\documentclass[a4paper,10pt]{report}\input{../0.Préambule/Préambule_cours_prof}\begin{document}

\pagestyle{empty}
\newpage
\section*{Enchainements d'opérations - Calcul littéral 1\hspace{3cm} Sujet A}
\addcontentsline{toc}{section}{Enchainements d'opérations - Calcul littéral 1\hspace{3cm} Sujet A}

\noindent \textit{L'usage de la calculatrice est \textbf{interdite}.} \hfill \textit{Durée : 55 minutes}

\paragraph{Exercice 1}: \quad Qu'est ce qu'une expression littérale ? % Que veut dire réduire une expression littérale ? 
\hfill \textbf{2 points}


\paragraph{Exercice 2}: \hspace{8cm} \textbf{5 points}

\medskip
\noindent Voici un programme de calcul 

\begin{enumerate}
\item Applique ce programme à chacun des nombres en écrivant une seule \\ expression permettant de trouver le résultat puis calcule :

\begin{enumerate}
\begin{tabularx}{\linewidth}{*{4}{>{\arraybackslash}X}}
\item $6$ &
\item $-6$ &&\\
\end{tabularx}
\end{enumerate}
\item Applique ce programme en choisissant la lettre $a$ au début du programme \\et écrit une seule expression permettant de trouver le résultat.
\end{enumerate}
\vspace{-46mm}
\begin{flushright}
\arrayrulecolor{black}
\renewcommand{\arraystretch}{0.1}
\begin{tabular}{|p{4.5cm}|}
\hline
\begin{verbatim} Choisis un nombre \end{verbatim}
\begin{verbatim} Calcule son carre \end{verbatim}
\begin{verbatim} Multiplie par (-2) \end{verbatim}
\begin{verbatim} Prendre l'opposé \end{verbatim}
\begin{verbatim} Soustraire (-5) \end{verbatim}
\\\hline
\end{tabular}
\end{flushright}

\paragraph{Exercice 3}: \hfill \textbf{3 points}

\medskip
\noindent Le premier mai, Ludo est allé vendre du muguet. Avec les $739$ brins cueillis, il avait composé $30$ gros bouquets de $12$ brins, des petits bouquets de $5$ brins et avait offert ses $4$ derniers brins de muguet à sa mère. %734 et 9

\medskip
\noindent Écris une expression qui permet de calculer le nombre de petits bouquets de Ludo puis calcule-la.

\paragraph{Exercice 4}: \hfill \textbf{5 points}

\medskip
\noindent Calcule \textbf{en détaillant} les étapes.

\begin{enumerate}
\begin{tabularx}{\linewidth}{*{3}{>{\arraybackslash}X}}
\item $I = (4 - 7) \times  5 - 2$&
%\item $J = 6 \div (-2) \times (8 - 2)$&
\item $K = 5 - 15 \div (10 - 5)$&
%\item $L = (31 - 13) \div 3 \times 2$&
\item $N = 10 - 5 \times \big(10 \div (-5) \big)$\\
%\item $M = 26 - \big( (-6) \times 5 + 36\big)$\\

\item $B = 6 \times \big[ (-2) \times (5 - 2) \big] $&
%\item $C = \big[(8 - 2) \times 8 - 4 \big] \div 4 - 18$&
\item $D = 1 - \big[ (31 - 5) \div 2 + 23) \big] \div 6 \div 2$&\\
%\item $E = 3 - \big[ 9 \times (8 \div 2) \big] \div 6 \times 3 - 2$&\\
\end{tabularx}
\end{enumerate}

\paragraph{Exercice 5}: \hfill \textbf{2 points}

\medskip
\noindent Supprime les parenthèses.

\begin{enumerate}
\begin{tabularx}{\linewidth}{*{4}{>{\arraybackslash}X}}
\item $-5+(-a+2b)$ &
%\item $9-(-x+3y)$ &
\item $-a-(-5+b-c)$ &
%\item $9+(-x^2-x+y)$ &
\item $7a+(-5+2b-c)$ &
%\item $-5-(x^2+4x)$ &
\item $b-(-a+6c-5)$ \\
%\item $-3z+(x-y-2)$\\
\end{tabularx}
\end{enumerate}

\paragraph{Exercice 6}: \hfill \textbf{3 points}

\medskip
\noindent Supprime les parenthèses puis réduis les expressions suivantes :

\begin{enumerate}
\begin{tabularx}{\linewidth}{*{3}{>{\arraybackslash}X}}
\item $L=5-(2x+3$) &
%\item $O=5x-(3-4x)$ &
\item $U=(x-4)-6$ &
%\item $P=(4x+2)+(-6x-2)$ &
\item $E=-(-3x-1)+(x-3)$ \\
%\item $R=8x-(5x+2)+(3-4x)$\\
\end{tabularx}
\end{enumerate}

\newpage
\section*{Enchainements d'opérations - Calcul littéral 1\hspace{3cm} Sujet B}
\addcontentsline{toc}{section}{Enchainements d'opérations - Calcul littéral 1\hspace{3cm} Sujet B}

\noindent \textit{L'usage de la calculatrice est \textbf{interdite}.} \hfill \textit{Durée : 55 minutes}

\paragraph{Exercice 1}: \quad Que veut dire réduire une expression littérale ? %Qu'est ce qu'une expression littérale ?   
\hfill \textbf{2 points}


\paragraph{Exercice 2}: \hspace{8cm} \textbf{5 points}

\medskip
\noindent Voici un programme de calcul 

\begin{enumerate}
\item Applique ce programme à chacun des nombres en écrivant une seule \\ expression permettant de trouver le résultat puis calcule :

\begin{enumerate}
\begin{tabularx}{\linewidth}{*{4}{>{\arraybackslash}X}}
\item $4$ &
\item $-4$ &&\\
\end{tabularx}
\end{enumerate}
\item Applique ce programme en choisissant la lettre $b$ au début du programme \\et écrit une seule expression permettant de trouver le résultat.
\end{enumerate}
\vspace{-46mm}
\begin{flushright}
\arrayrulecolor{black}
\renewcommand{\arraystretch}{0.1}
\begin{tabular}{|p{4.5cm}|}
\hline
\begin{verbatim} Choisis un nombre \end{verbatim}
\begin{verbatim} Calcule son carre \end{verbatim}
\begin{verbatim} Multiplie par (-2) \end{verbatim}
\begin{verbatim} Prendre l'opposé \end{verbatim}
\begin{verbatim} Soustraire (-6) \end{verbatim}
\\\hline
\end{tabular}
\end{flushright}

\paragraph{Exercice 3}: \hfill \textbf{3 points}

\medskip
\noindent Le premier mai, Ludo est allé vendre du muguet. Avec les $734$ brins cueillis, il avait composé $30$ gros bouquets de $12$ brins, des petits bouquets de $5$ brins et avait offert ses $9$ derniers brins de muguet à sa mère.

\medskip
\noindent Écris une expression qui permet de calculer le nombre de petits bouquets de Ludo puis calcule-la.

\paragraph{Exercice 4}: \hfill \textbf{5 points}

\medskip
\noindent Calcule \textbf{en détaillant} les étapes.

\begin{enumerate}
\begin{tabularx}{\linewidth}{*{3}{>{\arraybackslash}X}}
%\item $I = (4 - 7) \times  5 - 2$&
\item $J = 6 \div (-2) \times (8 - 2)$&
%\item $K = 5 - 15 \div (10 - 5)$&
\item $L = (31 - 13) \div 3 \times 2$&
%\item $N = 10 - 5 \times \big(10 \div (-5) \big)$\\
\item $M = 26 - \big( (-6) \times 5 + 36\big)$\\

%\item $B = 6 \times \big[ (-2) \times (5 - 2) \big] $&
\item $C = \big[(8 - 2) \times 8 - 4 \big] \div 4 - 18$&
%\item $D = 1 - \big[ (31 - 5) \div 2 + 23) \big] \div 6 \div 2$&\\
\item $E = 3 - \big[ 9 \times (8 \div 2) \big] \div 6 \times 3 - 2$&\\
\end{tabularx}
\end{enumerate}

\paragraph{Exercice 5}: \hfill \textbf{2 points}

\medskip
\noindent Supprime les parenthèses.

\begin{enumerate}
\begin{tabularx}{\linewidth}{*{4}{>{\arraybackslash}X}}
%\item $-5+(-a+2b)$ &
\item $9-(-x+3y)$ &
%\item $-a-(-5+b-c)$ &
\item $9+(-3-x+y)$ &
%\item $7a+(-5+2b-c)$ &
\item $-5-(u+4x)$ &
%\item $b-(-a+6c-5)$ \\
\item $-3z+(x-y-2)$\\
\end{tabularx}
\end{enumerate}

\paragraph{Exercice 6}: \hfill \textbf{3 points}

\medskip
\noindent Supprime les parenthèses puis réduis les expressions suivantes :

\begin{enumerate}
\begin{tabularx}{\linewidth}{*{3}{>{\arraybackslash}X}}
%\item $L=5+(2x+3$) &
\item $O=5-(3-4x)$ &
%\item $U=(x-4)-6$ &
\item $P=(4x+2)+(-6x-2)$ &
%\item $E=-(-3x-1)+(x-3)$ \\
\item $R=8x-(5x+2)$\\
\end{tabularx}
\end{enumerate}
\end{document}
